% GENERATED FILE. Edit canonical lessons, then run npm run generate:books.
% canonical-source-hash: fnv1a64:d449d0134f263701
% canonical-lessons: ES-C520-delito, ES-C520-crimen, ES-C520-delincuente, ES-C520-victima, ES-C520-investigacion

\chapter{El delito}
\label{ch:es-el-delito}
\hlchaptermodality{520}

\begin{chapteropening}
\textbf{By the end of this chapter:} \emph{I can talk about an offence and an offender, a crime, the victim and the investigation.}

Complete the last lesson of chapter 520: I can talk about an offence and an offender, a crime, the victim and the investigation.
\end{chapteropening}

\section[el delito]{el delito — the offence, the crime}

\label{lesson:ES-C520-delito}



\begin{quote}
Before the new one: say the Spanish for the patrol, then the Spanish for to chase.
\end{quote}

\subsection*{You'll want to know: el delito}
\textbf{el delito} — \textquotedblleft{}the offence, the crime\textquotedblright{}.

\begin{quote}
Robar es un delito.
\end{quote}

\begin{quote}
\emph{Stealing is a crime.}
\end{quote}

\emph{Cometer un delito} — to commit an offence.

\begin{cousinweb}
\emph{delito} is a learned borrowing from Latin \emph{dēlictum} 'a fault, a misdeed' (EtymDB), from \emph{dēlinquere} 'to fall short, to do wrong' (Lewis \& Short).
\end{cousinweb}

\begin{grammarlens}[title={What you've built}]
One more word for telling, explaining and describing at B1.
\end{grammarlens}

\subsection*{Your turn}
\begin{itemize}
\raggedright
  \item \emph{Say it:} \emph{el delito}
  \item \emph{Say it:} \emph{el delito}, once more
  \item \emph{Say it:} a sentence of your own with \emph{el delito}
  \item \emph{Recall:} say the Spanish for the patrol, then the Spanish for to chase, then say \emph{el delito} again
\end{itemize}

\begin{tcolorbox}[breakable,colback=teal!4,colframe=teal!35!black,title={Before you move on}]
What does el delito mean? (\textquotedblleft{}The offence, the crime\textquotedblright{}.) Say it once more.
\end{tcolorbox}
\section[el crimen]{el crimen — the crime (serious)}

\label{lesson:ES-C520-crimen}



\begin{quote}
Before the new one: say the Spanish for to chase, then the Spanish for the offence.
\end{quote}

\subsection*{You'll want to know: el crimen}
\textbf{el crimen} — \textquotedblleft{}the crime (serious)\textquotedblright{}.

\begin{quote}
Nadie vio el crimen.
\end{quote}

\begin{quote}
\emph{Nobody saw the crime.}
\end{quote}

The plural moves the accent: \emph{los crímenes}.

\begin{cousinweb}
\emph{crimen} is a learned borrowing from Latin \emph{crīmen} 'a verdict, a charge, a crime' (EtymDB).
\end{cousinweb}

\begin{grammarlens}[title={What you've built}]
One more word for telling, explaining and describing at B1.
\end{grammarlens}

\subsection*{Your turn}
\begin{itemize}
\raggedright
  \item \emph{Say it:} \emph{el crimen}
  \item \emph{Say it:} \emph{el crimen}, once more
  \item \emph{Say it:} a sentence of your own with \emph{el crimen}
  \item \emph{Recall:} say the Spanish for to chase, then the Spanish for the offence, then say \emph{el crimen} again
\end{itemize}

\begin{tcolorbox}[breakable,colback=teal!4,colframe=teal!35!black,title={Before you move on}]
What does el crimen mean? (\textquotedblleft{}The crime (serious)\textquotedblright{}.) Say it once more.
\end{tcolorbox}
\section[el delincuente]{el delincuente — the offender, the criminal}

\label{lesson:ES-C520-delincuente}



\begin{quote}
Before the new one: say the Spanish for the offence, then the Spanish for the crime.
\end{quote}

\subsection*{You'll want to know: el delincuente}
\textbf{el delincuente} — \textquotedblleft{}the offender, the criminal\textquotedblright{}.

\begin{quote}
Detuvieron al delincuente en la estación.
\end{quote}

\begin{quote}
\emph{They arrested the offender at the station.}
\end{quote}

The same form for a woman: \emph{la delincuente}.

\begin{cousinweb}
\emph{delincuente} is a learned borrowing from Latin \emph{dēlinquēns} 'doing wrong' (EtymDB), the participle of \emph{dēlinquere} 'to fall short, to do wrong'.
\end{cousinweb}

\begin{grammarlens}[title={What you've built}]
One more word for telling, explaining and describing at B1.
\end{grammarlens}

\subsection*{Your turn}
\begin{itemize}
\raggedright
  \item \emph{Say it:} \emph{el delincuente}
  \item \emph{Say it:} \emph{el delincuente}, once more
  \item \emph{Say it:} a sentence of your own with \emph{el delincuente}
  \item \emph{Recall:} say the Spanish for the offence, then the Spanish for the crime, then say \emph{el delincuente} again
\end{itemize}

\begin{tcolorbox}[breakable,colback=teal!4,colframe=teal!35!black,title={Before you move on}]
What does el delincuente mean? (\textquotedblleft{}The offender, the criminal\textquotedblright{}.) Say it once more.
\end{tcolorbox}
\section[la víctima]{la víctima — the victim}

\label{lesson:ES-C520-victima}



\begin{quote}
Before the new one: say the Spanish for the crime, then the Spanish for the offender.
\end{quote}

\subsection*{You'll want to know: la víctima}
\textbf{la víctima} — \textquotedblleft{}the victim\textquotedblright{}.

\begin{quote}
La víctima era un hombre de cuarenta años.
\end{quote}

\begin{quote}
\emph{The victim was a man of forty.}
\end{quote}

Always feminine, even for a man: \emph{la víctima}.

\begin{cousinweb}
\emph{víctima} is a learned borrowing from Latin \emph{victima} 'an animal offered in sacrifice' (EtymDB). The word moved from the altar to anyone who suffers harm.
\end{cousinweb}

\begin{grammarlens}[title={What you've built}]
One more word for telling, explaining and describing at B1.
\end{grammarlens}

\subsection*{Your turn}
\begin{itemize}
\raggedright
  \item \emph{Say it:} \emph{la víctima}
  \item \emph{Say it:} \emph{la víctima}, once more
  \item \emph{Say it:} a sentence of your own with \emph{la víctima}
  \item \emph{Recall:} say the Spanish for the crime, then the Spanish for the offender, then say \emph{la víctima} again
\end{itemize}

\begin{tcolorbox}[breakable,colback=teal!4,colframe=teal!35!black,title={Before you move on}]
What does la víctima mean? (\textquotedblleft{}The victim\textquotedblright{}.) Say it once more.
\end{tcolorbox}
\section[la investigación]{la investigación — the investigation; research}

\label{lesson:ES-C520-investigacion}



\begin{quote}
Before the new one: say the Spanish for the offender, then the Spanish for the victim.
\end{quote}

\subsection*{You'll want to know: la investigación}
\textbf{la investigación} — \textquotedblleft{}the investigation; research\textquotedblright{}.

\begin{quote}
La investigación duró dos años.
\end{quote}

\begin{quote}
\emph{The investigation lasted two years.}
\end{quote}

It also means scientific research: \emph{la investigación médica}.

\begin{cousinweb}
\emph{investigación} is a learned borrowing from Latin \emph{investīgātiō} 'a searching out', from \emph{investīgāre} 'to track, as a dog does' (Lewis \& Short), built on \emph{vestīgium} 'a footprint'.
\end{cousinweb}

\begin{grammarlens}[title={What you've built}]
One more word for telling, explaining and describing at B1.
\end{grammarlens}

\subsection*{Your turn}
\begin{itemize}
\raggedright
  \item \emph{Say it:} \emph{la investigación}
  \item \emph{Say it:} \emph{la investigación}, once more
  \item \emph{Say it:} a sentence of your own with \emph{la investigación}
  \item \emph{Recall:} say the Spanish for the offender, then the Spanish for the victim, then say \emph{la investigación} again
\end{itemize}

\begin{tcolorbox}[breakable,colback=teal!4,colframe=teal!35!black,title={Before you move on}]
What does la investigación mean? (\textquotedblleft{}The investigation; research\textquotedblright{}.) Say it once more.
\end{tcolorbox}
